\section{Requerimientos del Sistema}

La especificación de requerimientos constituye el contrato técnico que delimita
el alcance del sistema y establece los criterios verificables de aceptación. A
partir de los objetivos planteados y de las condiciones de operación del
parqueadero, los requerimientos se clasifican en funcionales, no funcionales y
de infraestructura.

\subsection{Requerimientos funcionales}

Los requerimientos funcionales describen las capacidades que el sistema debe
proporcionar para cumplir con su propósito operativo. La Tabla
\ref{tab:requerimientos_funcionales} los enumera y describe.

\begin{longtable}{
        P{0.07}
        P{0.35}
        P{0.50}
    }
    \caption{Requerimientos funcionales del sistema}
    \label{tab:requerimientos_funcionales}\\

    \toprule
    \rowcolor{gray!20}
    \textbf{ID} &
    \textbf{Requerimiento} &
    \textbf{Descripción} \\
    \midrule
    \endfirsthead

    \multicolumn{3}{c}{\textit{Continuación de la Tabla
    \ref{tab:requerimientos_funcionales}}} \\
    \toprule
    \rowcolor{gray!20}
    \textbf{ID} &
    \textbf{Requerimiento} &
    \textbf{Descripción} \\
    \midrule
    \endhead

    RF-01 & Adquisición de video &
    Recibir y procesar el flujo de video proveniente de las cámaras ubicadas en
    los puntos de acceso del parqueadero, con capacidad de operación continua.
    \\

    \midrule
    RF-02 & Detección de vehículos &
    Detectar y localizar cada vehículo presente en el área de aproximación,
    representándolo mediante una caja delimitadora.
    \\

    \midrule
    RF-03 & Identificación del tipo de vehículo &
    Clasificar el vehículo detectado en las categorías definidas para el entorno
    de estudio (por ejemplo, sedán, camioneta, minibús o motocicleta).
    \\

    \midrule
    RF-04 & Identificación del color predominante &
    Determinar el color predominante de la carrocería del vehículo dentro de un
    conjunto de clases cromáticas previamente definido.
    \\

    \midrule
    RF-05 & Detección de la placa de matrícula &
    Localizar la región de la placa de matrícula dentro de la caja del vehículo
    mediante una caja delimitadora específica.
    \\

    \midrule
    RF-06 & Reconocimiento de la matrícula &
    Reconocer los caracteres alfanuméricos de la placa detectada y construir la
    cadena de texto correspondiente al formato de matrícula vigente.
    \\

    \midrule
    RF-07 & Seguimiento del vehículo &
    Asignar y mantener un identificador temporal único a cada vehículo mientras
    permanece en el campo de visión de la cámara.
    \\

    \midrule
    RF-08 & Determinación de eventos &
    Clasificar el movimiento del vehículo como ingreso o salida a partir del
    cruce de una línea virtual de control y del sentido de su trayectoria.
    \\

    \midrule
    RF-09 & Registro del evento &
    Almacenar cada evento con la placa reconocida, el tipo y color del vehículo,
    la fecha, la hora, la cámara y el sentido correspondiente.
    \\

    \midrule
    RF-10 & Control de autorización &
    Verificar la matrícula reconocida contra los registros de vehículos
    autorizados y reportar los accesos no autorizados.
    \\

    \midrule
    RF-11 & Consulta de información &
    Permitir la consulta y visualización del historial de eventos vehiculares
    registrados por el sistema.
    \\

    \midrule
    RF-12 & Generación de estadísticas &
    Producir información estadística a partir de los registros históricos, como
    afluencia vehicular y tiempos de permanencia.
    \\

    \midrule
    RF-13 & Gestión de usuarios &
    Administrar el acceso de los usuarios al sistema web conforme a roles
    definidos (administrador y operador).
    \\*
    \multicolumn{3}{r}{\footnotesize\textit{Fuente: Elaboración propia}} \\*
    \bottomrule
\end{longtable}

El requerimiento RF-10 se conserva en la especificación por trazabilidad con el
perfil del proyecto, pero se declara \textbf{fuera del alcance implementado}: el
contraste contra un padrón de vehículos habilitados exige un registro previo que
esta entrega no construye, de modo que el sistema documenta el acceso sin
clasificarlo como autorizado. La limitación correspondiente se declara en el
capítulo introductorio y el estado del componente, en el capítulo de
implementación.

\subsection{Requerimientos no funcionales}

Los requerimientos no funcionales establecen restricciones de calidad y
condiciones de operación. Su cumplimiento se verifica en la fase de evaluación
mediante las métricas descritas en el marco teórico; los valores objetivo son
coherentes con las hipótesis planteadas en el perfil del proyecto.

\begin{longtable}{
        P{0.09}
        P{0.28}
        P{0.55}
    }
    \caption{Requerimientos no funcionales del sistema}
    \label{tab:requerimientos_no_funcionales}\\

    \toprule
    \rowcolor{gray!20}
    \textbf{ID} &
    \textbf{Categoría} &
    \textbf{Criterio de cumplimiento} \\
    \midrule
    \endfirsthead

    \multicolumn{3}{c}{\textit{Continuación de la Tabla
    \ref{tab:requerimientos_no_funcionales}}} \\
    \toprule
    \rowcolor{gray!20}
    \textbf{ID} &
    \textbf{Categoría} &
    \textbf{Criterio de cumplimiento} \\
    \midrule
    \endhead

    RNF-01 & Rendimiento temporal &
    El \textit{pipeline} completo de visión debe procesar un evento de acceso con
    una latencia promedio igual o inferior a un segundo.
    \\

    \midrule
    RNF-02 & Exactitud de detección &
    El modelo de detección de vehículos debe alcanzar un mAP@50 igual o superior
    al 80\,\%.
    \\

    \midrule
    RNF-03 & Exactitud de clasificación &
    Los modelos de identificación de tipo y color predominante deben alcanzar un
    F1-score macro igual o superior al 85\,\%.
    \\

    \midrule
    RNF-04 & Exactitud de reconocimiento &
    El módulo de reconocimiento de placas debe alcanzar un Recall igual o
    superior al 90\,\% en la detección y una exactitud por carácter igual o
    superior al 90\,\% en la lectura.
    \\

    \midrule
    RNF-05 & Exactitud de eventos &
    La determinación de los eventos de ingreso y salida debe alcanzar una
    exactitud igual o superior al 90\,\%.
    \\

    \midrule
    RNF-06 & Disponibilidad &
    El sistema debe operar de forma continua durante la jornada de operación del
    parqueadero, resistiendo reinicios controlados de sus componentes.
    \\

    \midrule
    RNF-07 & Escalabilidad &
    La arquitectura debe permitir incorporar nuevos puntos de acceso sin
    modificar el núcleo del subsistema de visión ni degradar el rendimiento.
    \\

    \midrule
    RNF-08 & Eficiencia en el borde &
    El motor de percepción debe ejecutarse en hardware con recursos
    restringidos, empleando modelos optimizados y cuantizados.
    \\

    \midrule
    RNF-09 & Seguridad y privacidad &
    El acceso a la información y a las imágenes debe estar protegido, y el
    tratamiento de datos debe considerar las disposiciones aplicables sobre
    protección de datos personales.
    \\

    \midrule
    RNF-10 & Trazabilidad &
    El sistema debe registrar eventos y errores en bitácoras que permitan
    auditar su comportamiento y reconstruir un evento ante una incidencia.
    \\*
    \multicolumn{3}{r}{\footnotesize\textit{Fuente: Elaboración propia}} \\*
    \bottomrule
\end{longtable}

\subsection{Requerimientos de infraestructura y entorno}

El sistema se despliega sobre una infraestructura física mínima compuesta por
dispositivos de captura en los puntos de acceso, una unidad de procesamiento en
el borde encargada de la inferencia, un servidor para la plataforma de gestión y
la base de datos, y la red local que interconecta estos componentes. La Tabla
\ref{tab:infraestructura} resume la función esperada de cada componente.

\begin{longtable}{
        P{0.22}
        P{0.35}
        P{0.35}
    }
    \caption{Componentes de infraestructura y entorno}
    \label{tab:infraestructura}\\

    \toprule
    \rowcolor{gray!20}
    \textbf{Componente} &
    \textbf{Función en el sistema} &
    \textbf{Consideración de diseño} \\
    \midrule
    \endfirsthead

    \multicolumn{3}{c}{\textit{Continuación de la Tabla
    \ref{tab:infraestructura}}} \\
    \toprule
    \rowcolor{gray!20}
    \textbf{Componente} &
    \textbf{Función en el sistema} &
    \textbf{Consideración de diseño} \\
    \midrule
    \endhead

    Cámara del punto de acceso &
    Capturar las imágenes y el video de los vehículos que aproximan a la
    barrera. &
    La posición, altura y ángulo de montaje determinan la calidad de la placa en
    el encuadre.
    \\

    \midrule
    Unidad de procesamiento en el borde &
    Ejecutar los modelos de detección, reconocimiento y seguimiento. &
    Recursos limitados de cómputo y memoria; requiere modelos cuantizados.
    \\

    \midrule
    Servidor de aplicación y base de datos &
    Alojar la API de gestión, la lógica de autorización y la persistencia de los
    eventos. &
    Debe garantizar disponibilidad durante la operación y respaldo de la
    información.
    \\

    \midrule
    Red local &
    Transportar el video desde la cámara y los eventos hacia el servidor. &
    El ancho de banda y la estabilidad condicionan la latencia del sistema.
    \\

    \midrule
    Almacenamiento &
    Conservar los registros de eventos y, cuando corresponda, la evidencia
    visual asociada. &
    La política de retención debe equilibrar trazabilidad y privacidad.
    \\*
    \multicolumn{3}{r}{\footnotesize\textit{Fuente: Elaboración propia}} \\*
    \bottomrule
\end{longtable}

Las características técnicas concretas de los dispositivos (modelo de cámara,
capacidad de cómputo del equipo de borde y dimensionamiento del servidor) se
determinan en la fase de análisis de las condiciones del parqueadero y se
documentan en el capítulo de implementación.
